% !TeX program = xelatex
% 使用 XeLaTeX 编译，需要中文字体支持（如 Windows 下的 SimHei/SimSun，或 macOS 下的 Heiti SC/Songti SC）
\documentclass[zihao=-4, linespread=1.25]{ctexart} % 小四号字，1.25倍行距

% 页面设置（A4，上下左右边距可按需调整）
\usepackage[top=2.54cm, bottom=2.54cm, left=3.17cm, right=3.17cm]{geometry}

% 必要宏包
\usepackage{amsmath}          % 数学公式
\usepackage{graphicx}         % 插图
\usepackage{caption}          % 图题、表题定制
\usepackage{booktabs}         % 三线表
\usepackage{enumitem}         % 列表定制（如需要）
\usepackage{fontspec}         % 字体设置（配合 xelatex）
\usepackage{titlesec}         % 标题格式
\usepackage{titletoc}         % 目录格式（可选）
\usepackage{setspace}         % 行距
\usepackage{indentfirst}      % 首行缩进
\usepackage[colorlinks, citecolor=blue]{hyperref} % 参考文献超链接
\usepackage{chngcntr}          % 提供 \counterwithin 命令
\usepackage{tabularx}
\usepackage[super,sort&compress]{natbib}

% ========== 中文字体设置（根据操作系统修改） ==========
\setCJKmainfont{SimSun}[AutoFakeBold=2.5]
\setCJKsansfont{SimSun}[AutoFakeBold=2.5]
% macOS 示例（取消下面注释，并注释上面两行）：
% \setCJKmainfont{Songti SC}[AutoFakeBold=2.5]
% \setCJKsansfont{Heiti SC}[AutoFakeBold=2.5]

% 西文字体
\setmainfont{Times New Roman}

% ========== 全局段落格式 ==========
\setlength{\parindent}{2em}      % 首行缩进两格（中文字符两个）
\setlength{\parskip}{0pt}        % 段落间不额外空行

% ========== 标题格式 ==========
% 一级标题（黑体小二号居中，段前距12磅≈4.23mm）
\titleformat{\section}
  {\centering\sffamily\bfseries\zihao{-2}}  % 小二号，黑体，居中
  {\thesection}{1em}{}
\titlespacing{\section}{0pt}{12bp}{0bp}    % 段前12磅

% 二级标题（黑体四号左对齐，段前距12磅）
\titleformat{\subsection}
  {\sffamily\bfseries\zihao{4}}             % 四号，黑体，左对齐
  {\thesubsection}{1em}{}
\titlespacing{\subsection}{0pt}{12bp}{0bp}  % 段前12磅

% 三级标题（黑体小四号左对齐，段前距12磅）
\titleformat{\subsubsection}
  {\sffamily\bfseries\zihao{-4}}            % 小四号，黑体
  {\thesubsubsection}{1em}{}
\titlespacing{\subsubsection}{0pt}{12bp}{0bp}

% ========== 图表标题格式 ==========
\captionsetup[table]{
  labelsep=space,           % 编号与标题之间用空格
  font={bf,small},          % 五号粗体（small 对应五号）
  position=above,           % 表题在表格上方
  skip=2pt
}
\captionsetup[figure]{
  labelsep=space,
  font={bf,small},
  position=below,           % 图题在图片下方
  skip=2pt
}
\counterwithin{table}{section} % 表格编号包含章节号
\numberwithin{equation}{section} % 公式编号包含章节号
\renewcommand{\theequation}{\thesection-\arabic{equation}}
\numberwithin{figure}{section} % 图片编号包含章节号
\renewcommand{\thefigure}{\thesection-\arabic{figure}}

% 重定义表格标题命令，使其符合“表1.1 表名”且不加冒号
\renewcommand{\tablename}{表}
\renewcommand{\figurename}{图}

% ========== 摘要、关键词环境 ==========
\newenvironment{abstractcn}{%
  \par\noindent{\sffamily\bfseries\zihao{-4} 摘要：}\ignorespaces
}{%
  \par\vspace{6bp}
}
\newenvironment{keywordscn}{%
  \par\noindent{\sffamily\bfseries\zihao{-4} 关键词：}\ignorespaces
}{%
  \par\vspace{12bp}
}

% ========== 参考文献格式 ==========
\bibliographystyle{plain}   % 可根据需要改为 gbt7714 等
\renewcommand{\refname}{参考文献}  % 将默认的 Bibliography 改为“参考文献”

% 自定义参考文献条目格式（首行缩进两格，序号和内容空一格）
\makeatletter
\renewenvironment{thebibliography}[1]{%
  \section*{\centering\sffamily\bfseries\zihao{-2} 参考文献}% 使用一级标题样式
  \list{\@biblabel{\@arabic\c@enumiv}}%
       {\settowidth\labelwidth{\@biblabel{#1}}%
        \leftmargin\labelwidth
        \advance\leftmargin\labelsep
        \@openbib@code
        \usecounter{enumiv}%
        \let\p@enumiv\@empty
        \renewcommand\theenumiv{\@arabic\c@enumiv}%
        \setlength{\itemindent}{2em}      % 首行缩进两格
        \setlength{\leftmargin}{2em}      % 序号缩进
        \setlength{\labelwidth}{1.5em}    % 序号宽度
        \setlength{\labelsep}{0.5em}      % 序号与内容间隔
        \setlength{\itemsep}{0pt}         % 条目间距
        \fontsize{12pt}{18pt}\selectfont  % 小四号（12pt）
        }
  \sloppy
  \clubpenalty4000
  \@clubpenalty \clubpenalty
  \widowpenalty4000%
  \sfcode`\.\@m}
{\def\@noitemerr
  {\@latex@warning{Empty `thebibliography' environment}}%
  \endlist}
\makeatother

\usepackage{fancyhdr}
\pagestyle{fancy}
\fancyhf{}   % 清空所有页眉页脚

% 设置页眉左侧为主标题，页眉右侧为页码
\fancyhead[L]{\zihao{5}\songti 标题（开题报告）（宋体五号字）}
\fancyhead[R]{\zihao{5}\songti \thepage}

% 删除默认的页眉横线（可选，若需要横线则保留下一行）
\renewcommand{\headrulewidth}{0.4pt}  % 横线宽度，不需要可设为0pt

% ========== 文档开始 ==========
\begin{document}

% ---------- 标题部分 ----------
\begin{center}
  \vspace{12bp}  % 段前距12磅
  {\sffamily\bfseries\zihao{-2} 标题（黑体小二号粗体并居中，段前距为12磅）}  % 标题小二号黑体居中
  \vspace{12bp}

  {\sffamily\bfseries\zihao{4} 班级\ 姓{\quad}名（学号）（班级后在半角状态下空一格，黑体四号并居中，段后距为12磅）}  % 班级 姓名（学号）
  \vspace{12bp}
\end{center}

% ---------- 摘要 ----------
\begin{abstractcn}
  具体内容（宋体小四号，超出一行顶格排版）
\end{abstractcn}

% ---------- 关键词 ----------
\begin{keywordscn}
  关键词1，关键词2，……（宋体小四号，3～5词，中间用逗号，结束不用标点，超出一行顶格排版）
\end{keywordscn}

% ---------- 正文 ----------
\vspace{12bp}

\section{研究背景（黑体小二号居中，段前距为12磅）}

内容（首行缩进两格，宋体小四号）

\subsection{二级标题（黑体四号，左对齐，段前距为12磅）}

内容（首行缩进两格，宋体小四号）

\subsubsection{三级标题（黑体小四号，左对齐，段前距为12磅）}

内容（首行缩进两格，宋体小四号）

\begin{table}[htbp]
  \centering
  {\zihao{5}\songti                 % 开始宋体五号
  \caption{表名（编号与表名相距半角状态下的一个空格，表正上方，宋体五号粗体）}
  \label{tab:attention}
  \begin{tabularx}{\textwidth}{XXX}
    \toprule
    模块名称 & 参数量（M） & mIoU（\%） \\
    \midrule
    SE       & 0.21 & 85.3 \\
    CBAM     & 0.35 & 86.1 \\
    Ours     & 0.28 & 87.6 \\
    \bottomrule
  \end{tabularx}
  }
\end{table}

\subsection{二级标题（黑体四号，左对齐，段前距为12磅）}

\subsubsection{三级标题（黑体小四号，左对齐，段前距为12磅）}

内容（首行缩进两格，宋体小四号）

\section{文献综述（黑体小二号居中，段前距为12磅）}

\subsection{二级标题（黑体四号，左对齐，段前距为12磅）}

\subsubsection{三级标题（黑体小四号，左对齐，段前距为12磅）}

内容\cite{ref1}（内容首行缩进两格，宋体小四号；引文在右上角标注，五号字）

\begin{equation}
  A = -\sum_{i} y_i \log(\hat{y}_i) \label{eq:loss}
\end{equation}

\subsection{二级标题（黑体四号，左对齐，段前距为12磅）}

\subsubsection{三级标题（黑体小四号，左对齐，段前距为12磅）}

内容\cite{ref2,ref3,ref4}（内容首行缩进两格，宋体小四号；引文在右上角标注，五号字）

\begin{figure}[htbp]
  \centering
  \includegraphics[width=0.6\textwidth]{example-image}  % 替换为实际图片
  \caption{图名（编号与图名相距半角状态下的一个空格，图正下方，宋体五号粗体）}
  \label{fig:unet}
\end{figure}

\section{技术路线（黑体小二号居中，段前距为12磅）}

内容（首行缩进两格，宋体小四号）

\section{进度安排（黑体小二号居中，段前距为12磅）}

内容（首行缩进两格，宋体小四号）

% ---------- 参考文献 ----------
\begin{thebibliography}{99}
\bibitem{ref1} 著录内容（首行缩进两格，小四号字，序号和著录内容之间在半角状态下空一格）
\bibitem{ref2} 著录内容（首行缩进两格，小四号字，序号和著录内容之间在半角状态下空一格）
\bibitem{ref3} 著录内容（首行缩进两格，小四号字，序号和著录内容之间在半角状态下空一格）
\bibitem{ref4} 著录内容（首行缩进两格，小四号字，序号和著录内容之间在半角状态下空一格）
\end{thebibliography}

\end{document}